% ============================================================
\chapter{Indexing and Data Structures}
\label{ch:indexing}
% ============================================================

Imagine a library of ten million books without a catalogue. To find a volume, one would need to walk through every shelf, one book at a time. This is exactly what a DBMS does during a \emph{full table scan}: it examines every row of the table. An \emph{index} is the database's catalogue: an auxiliary structure that locates relevant rows in $O(\log n)$ instead of $O(n)$. The B+ tree, invented by Rudolf Bayer and Edward McCreight in 1972, is the dominant indexing structure: balanced, efficient for equality and range searches, and optimized for disk access. Hash indexes, bitmap indexes, and spatial indexes (R-tree) complete the arsenal. This chapter develops these structures and the indexing strategies that make the difference between a query in milliseconds and one in minutes.

\begin{intuition}[title=\textbf{Central idea}]
An index is an auxiliary data structure that accelerates lookups in a table
by avoiding a full table scan. Choosing the right index type and the right
columns to index is essential for performance.
\end{intuition}

% ------------------------------------------------------------
\section{Motivation}
% ------------------------------------------------------------

Without an index, any lookup in a table of $n$ rows requires a sequential
scan in $O(n)$. A well-chosen index reduces the complexity to
$O(\log n)$ or even $O(1)$.

\begin{remark}
An index speeds up reads but slows down writes (inserts, updates, deletes)
because the structure must be maintained. The read/write trade-off guides
index selection.
\end{remark}

% ------------------------------------------------------------
\section{B-trees and B$^+$-trees}
% ------------------------------------------------------------

\begin{definition}[B-tree of order $m$]
\index{B-tree}
A \textbf{B-tree} of order $m$ is a balanced search tree where each
internal node has at most $m$ children and at least $\lceil m/2 \rceil$
children (except the root). Keys are sorted within each node and data
is stored in all nodes.
\end{definition}

\begin{definition}[B$^+$-tree]
\index{B+-tree}
A \textbf{B$^+$-tree} is a variant where:
\begin{itemize}
  \item data is stored only in the \textbf{leaves},
  \item internal nodes contain only routing keys,
  \item leaves are linked in a doubly linked list for sequential scans.
\end{itemize}
\end{definition}

\begin{proposition}[B$^+$-tree complexity]
For a B$^+$-tree of order $m$ containing $n$ keys:
\begin{itemize}
  \item Height: $h = O(\log_m n)$.
  \item Search, insertion, deletion: $O(\log_m n)$ disk accesses.
  \item Range scan of $k$ keys: $O(\log_m n + k/B)$ where $B$ is the
        number of keys per leaf.
\end{itemize}
\end{proposition}

\begin{attention}
The B$^+$-tree is the default indexing structure in most relational DBMS
(PostgreSQL, MySQL/InnoDB, Oracle, SQL Server). Do not confuse it with a
binary search tree (BST), which is not suited for disk-based access.
\end{attention}

\begin{minted}{python}
-- Creating a B+ index (implicit) in SQL
CREATE INDEX idx_name ON students (last_name);

-- Composite index
CREATE INDEX idx_name_first ON students (last_name, first_name);

-- Verify index usage
EXPLAIN ANALYZE SELECT * FROM students WHERE last_name = 'Smith';
\end{minted}

% ------------------------------------------------------------
\section{Hash indexes}
% ------------------------------------------------------------

\begin{definition}[Hash index]
\index{hash index}
A \textbf{hash index} uses a function $h : \text{key} \to \{0, 1, \ldots,
B-1\}$ to map a key directly to a bucket. Equality lookups run in $O(1)$
on average.
\end{definition}

\begin{remark}
Hash indexes do \textbf{not} support range queries
(\texttt{WHERE x BETWEEN a AND b}) or ordering. They are suitable only
for exact equality lookups.
\end{remark}

\begin{definition}[Extendible hashing]
\index{extendible hashing}
\textbf{Extendible hashing} uses a directory of pointers indexed by the
first $d$ bits of the hash. When a bucket overflows, the directory is
doubled (or the bucket is split) without reorganizing the entire table.
\end{definition}

% ------------------------------------------------------------
\section{Bitmap indexes}
% ------------------------------------------------------------

\begin{definition}[Bitmap index]
\index{bitmap index}
A \textbf{bitmap index} associates each distinct value $v$ of a column
with a bit vector $B_v$ of length $n$ (number of rows):
$B_v[i] = 1$ if row $i$ has value $v$, 0 otherwise.
\end{definition}

\begin{proposition}[Efficiency of logical operations]
Queries combining multiple predicates translate into fast bitwise
operations:
\begin{itemize}
  \item \texttt{WHERE color = 'red' AND size = 'L'} $\to$ $B_{\text{red}} \wedge B_{\text{L}}$.
  \item \texttt{WHERE color = 'red' OR size = 'L'} $\to$ $B_{\text{red}} \vee B_{\text{L}}$.
\end{itemize}
Counting is done via \texttt{POPCOUNT} in $O(n/w)$ where $w$ is the
machine word size.
\end{proposition}

\begin{remark}
Bitmap indexes are particularly effective for \textbf{low-cardinality}
columns (gender, status, category). They are used extensively in data
warehouses (Oracle, Vertica).
\end{remark}

% ------------------------------------------------------------
\section{Spatial indexes: R-trees}
% ------------------------------------------------------------

\begin{definition}[R-tree]
\index{R-tree}
An \textbf{R-tree} is a balanced tree for indexing multidimensional spatial
data. Each node contains a \textbf{minimum bounding rectangle} (MBR) that
encloses all objects in its subtree.
\end{definition}

\begin{example}
Spatial search: ``Find all restaurants within 500m'':
\begin{minted}{python}
-- PostGIS (spatial extension for PostgreSQL)
CREATE INDEX idx_geom ON restaurants USING GIST (geom);

SELECT name, ST_Distance(geom, ST_MakePoint(-73.99, 40.75)::geography) AS dist
FROM restaurants
WHERE ST_DWithin(geom, ST_MakePoint(-73.99, 40.75)::geography, 500)
ORDER BY dist;
\end{minted}
\end{example}

% ------------------------------------------------------------
\section{Full-text indexes}
% ------------------------------------------------------------

\begin{definition}[Inverted index]
\index{inverted index}
An \textbf{inverted index} maps each term (word) to the list of documents
(rows) that contain it, optionally with position and frequency information.
\end{definition}

\begin{example}
\begin{minted}{python}
-- PostgreSQL: GIN index for full-text search
CREATE INDEX idx_fts ON articles USING GIN (to_tsvector('english', content));

SELECT title
FROM articles
WHERE to_tsvector('english', content) @@ to_tsquery('english', 'bayesian & statistics');
\end{minted}
\end{example}

% ------------------------------------------------------------
\section{Index selection strategies}
% ------------------------------------------------------------

\begin{definition}[Index selection problem]
\index{index selection}
Given a workload (set of queries with their frequencies) and a disk space
budget, choose the set of indexes that minimizes total execution cost.
\end{definition}

Heuristic rules:
\begin{enumerate}
  \item Index columns appearing in \texttt{WHERE}, \texttt{JOIN}, and
        \texttt{ORDER BY} clauses.
  \item Prefer covering composite indexes to avoid table access.
  \item Avoid indexing frequently modified high-cardinality columns.
  \item Use \texttt{EXPLAIN ANALYZE} to verify actual index usage.
\end{enumerate}

\begin{example}
\begin{minted}{python}
-- Covering index: all columns in the query are in the index
CREATE INDEX idx_covering ON orders (customer_id, order_date)
    INCLUDE (amount);

-- The following query can be satisfied without accessing the table
SELECT order_date, amount
FROM orders
WHERE customer_id = 42
ORDER BY order_date DESC;
\end{minted}
\end{example}

% ------------------------------------------------------------
\section{Key formulas}
% ------------------------------------------------------------

\begin{keyformulas}
\begin{itemize}
  \item \textbf{B$^+$-tree}: search in $O(\log_m n)$ disk accesses, range scan in $O(\log_m n + k/B)$
  \item \textbf{Hash}: equality lookup in $O(1)$, no range queries
  \item \textbf{Bitmap}: logical operations in $O(n/w)$, ideal for low cardinality
  \item \textbf{R-tree}: spatial search in $O(\log n)$ on average
  \item \textbf{Trade-off}: read acceleration vs.\ write overhead
\end{itemize}
\end{keyformulas}

% ------------------------------------------------------------
\section{Exercises}
% ------------------------------------------------------------

\begin{exercise}
Draw the B$^+$-tree of order 4 obtained after sequentially inserting keys
$3, 7, 1, 14, 8, 5, 11, 17, 13, 6, 23, 12, 20, 26, 4, 16$. Show the
state after deleting key 14.
\end{exercise}

\begin{exercise}
A table \texttt{employees} has 10 million rows. The column
\texttt{department} has 50 distinct values; the column \texttt{salary}
is continuous. What type of index do you recommend for each? Justify.
\end{exercise}

\begin{exercise}
Write a SQL query that \emph{cannot} use a hash index but \emph{can} use a
B$^+$-tree index. Explain why.
\end{exercise}

\begin{exercise}
We have a bitmap index on the column \texttt{status} (3 values: active,
inactive, suspended) and a bitmap index on \texttt{region} (10 values).
How many bitwise AND operations are needed to answer
\texttt{WHERE status = 'active' AND region IN ('NY', 'CA')}?
\end{exercise}

\begin{exercise}
Use \texttt{EXPLAIN ANALYZE} in PostgreSQL to compare the performance of a
query with and without an index. Measure the number of pages read (buffers)
and execution time.
\end{exercise}
